% ======================================================================
% experiments/dcache.tex — Level 2: enabling the D-cache
%
% ★ 軸はボードではなく「キャッシュ構成」。ゆえに boards/ ではなく
%    experiments/ に置く。同じ枠組みで別の軸を切りたくなったら
%    experiments/ に足す（例: 動的ワークスティーリング = Level 3）。
%
%    設計・制約・DMA棚卸し・ハーネスは測定に依存せず確定。
%    D-cache ON の「スケールするか」の実測値だけがシリアル取得待ち
%    （§末尾の \TODO を、数値が出たら差し替える）。
% ======================================================================

\subsection{動機：支配変数はキャッシュ構成である}

第 \ref{sec:rpi3} 節の測定で最も効いていた制約は、SMP の方式ではなく
\textbf{D-cache が OFF であること}であった。メモリ律速がスケールしなかったのも
（\bench{fill} \spd{0.93}）、\texttt{volatile} を外したら \spd{3.04} が
\spd{3.66} に上がったのも（第 \ref{sec:rpi3-volatile} 項）、動的ワークスティーリングを
断念して静的分割にしたのも、すべて D-cache OFF が原因である。

そこで問いは：\textbf{D-cache を ON にすると、メモリ律速の処理はスケールし始めるか}。
もしそうなら「ウィンドウシステムは並列化しない」という第 \ref{sec:rpi3}\,節の
結論が覆る。この問いは SMP の方式とは独立であり、キャッシュ構成という別の軸に属する。

\subsection{D-cache ON は「フラグを 1 にする」ではない}

素朴に \texttt{SCTLR.C} を 1 にするのは\textbf{危険で、過去に一度 Pi 3 を
brick させている}（\texttt{mmu.c} 冒頭に記録あり）。安全に ON にするには、
この移植に 1 行も存在しなかったキャッシュ管理を新規に書く必要があった。
本実験で顕在化した罠は 4 つある。

\begin{enumerate}[leftmargin=1.6em]
  \item \textbf{リセット直後のキャッシュはゼロでなくゴミ。}
        invalidate せずに \texttt{C=1} にすると、コアはゴミを有効データと
        見なす。これが brick の原因である。set/way invalidate ループを
        \texttt{C=1} の\textbf{前に}実行しなければならない（後では、いま
        キャッシュした本物のデータを捨ててしまう）。

  \item \textbf{Cortex-A53 は SMPEN なしでは非コヒーレント。}
        \texttt{mmu.c} は RAM を shareable にマップしているが、
        A53 では shareable かつ cacheable なメモリでも \texttt{CPUECTLR.SMPEN}
        を立てるまで\textbf{コア間でコヒーレントにならない}。\texttt{C=1} かつ
        SMPEN=0 だと、第 \ref{sec:design} 節のロックフリー・ジョブメールボックス
        （\texttt{volatile} + \texttt{dsb}）が静かに壊れる。\texttt{volatile} は
        コンパイラを縛るだけで、キャッシュは縛らない。

  \item \textbf{二次コアの起動窓。}
        二次コアは \texttt{\_smp\_start} の時点で\textbf{まだ MMU も
        キャッシュも OFF}であり、その状態でスタックポインタを RAM から読む。
        この読みはコア 0 のキャッシュをスヌープしない。ゆえにコア 0 は
        スタックアドレスを\textbf{起こす前に clean}せねばならず、さもなくば
        二次コアは null スタックで起きて最初の push で死ぬ。

  \item \textbf{\texttt{mmu\_disable} の前提が崩れる。}
        「キャッシュ OFF だから flush 不要」というコメントが \texttt{C=1} で
        偽になる。\texttt{kexec} 前にキャッシュを落とすと、\texttt{/upload} が
        キャッシュ経由で書いた次カーネルのイメージが dirty ライン内に取り残され、
        \textbf{stale メモリへジャンプ}する。clean を追加した。
\end{enumerate}

\noindent
これらを \texttt{cache.c}（set/way・範囲・SMPEN の各操作を新規実装）と、
\texttt{mmu.c} の \texttt{DCACHE\_ON} 経路（順序が命：invalidate と SMPEN を
\texttt{C=1} の前に）で処理した。

\subsection{最大の障害：USB DMA はページ属性で直せない}
\label{sec:dcache-usb}

D-cache ON でも DMA するハードとの共有メモリは非コヒーレントになる。
そこでカーネル全体の\textbf{DMA 地点を棚卸し}した。結果は明快で、
FB・電源・WiFi のメールボックスや SD/SDIO は\textbf{全て PIO}（DATA レジスタ
経由）であり、コヒーレンシ問題を持たない。唯一の DMA エンジンは
\textbf{DWC USB}であり、これが有線 LAN を担っている。

そして USB こそが\textbf{構造的に直せない}地点であった。DWC ドライバは
呼び出し側のポインタをそのまま DMA アドレスに渡し、呼び出し側は
\textbf{ヒープ}バッファ（\texttt{memget} は 8 バイト境界＝隣の割り当てと
64 バイトのキャッシュラインを共有）と、あろうことか\textbf{スタック上の
局所変数}（\texttt{smsc9512.c} のレジスタ操作が 4 バイトを DMA する）を渡す。
64 バイトのラインを invalidate すれば、\textbf{同じラインにある戻り番地や
退避レジスタごと破棄}する。これはページ属性では解けず、バウンスバッファと
ドライバ全域の 64 バイト境界化を要する。

この発見自体が結論である：\textbf{D-cache を入れられないのはキャッシュ管理を
書くのが面倒だからではなく、ドライバの DMA バッファの出自が構造的に許さない
から}である。

\subsection{実験の切り分け：USB を解かずに問いに答える}

問いは「D-cache でメモリ律速はスケールするか」であり、これに答えるのに
USB を解く必要はない。そこで実験用ビルド変種を用意した。

\begin{center}
\texttt{-DDCACHE\_EXPERIMENT -DDCACHE\_ON -DOS\_MINIMAL}
\end{center}

\noindent
\texttt{OS\_MINIMAL} でネットと WM を落とし、\texttt{DCACHE\_EXPERIMENT} で
さらに \texttt{usbinit}（唯一の DMA エンジン）を落とす。これで DMA する相手が
一つも存在しない状態で \texttt{C=1} にでき、USB の構造問題を回避したまま
本題を測れる。ただし\textbf{HDMI は残した}——ユーザ要求であり、また
フレームバッファの扱いが D-cache ON 実装の良い試金石でもあるためである。

HDMI を残すには 2 つの手当てが要った。FB のメールボックス要求構造体は
元は\textbf{スタック上}にあり（GPU が stale RAM を読み \texttt{address=0}
＝SYSERR ＝画面出ず）、64 バイト境界・64 バイト詰めの専用 static に移して
handshake の前後で clean/invalidate した。FB のピクセル領域自体は、
GPU がアドレスを返した\textbf{後に} \texttt{mmu\_set\_range\_noncached()} で
Normal Non-cacheable へ貼り替えた（アドレスは \texttt{mmu\_init} の後に
判明するため、静的な表には焼けない）。

\subsection{測定ハーネス}

実験変種はネットが無いので \route{/bench} を提供できない。そこでワークロード本体を
\texttt{apps/smpbench.c} に切り出し、\textbf{HTTP 経路（既定カーネル）と起動時
コンソール経路（実験カーネル）が同一のコードを呼ぶ}ようにした。キャッシュ
ON/OFF の比較は、両者が同一実装を走らせて初めて意味を持つ——別実装なら
「キャッシュ設定の差」ではなく「実装の差」を測ってしまう。この切り出しの
過程で、第 \ref{sec:rpi3-measbug} 項の測定バグを発見・修正している。

実験カーネルは起動時に \texttt{smpbench\_run\_all()} を実行し、\bench{fill} の
サイズ掃引を中心とした結果をシリアルへ出力する。

\subsection{実測結果}
\label{sec:dcache-result}

D-cache ON カーネルを実機に投入し、起動時ハーネスの出力をシリアルで取得した。
まず 3 つの前提が\textbf{すべて確認された}。

\begin{itemize}[leftmargin=1.6em]
  \item \textbf{起動した（brick 回避）}。Xinu のバナーからシェルまで正常に立ち
        上がった。set/way invalidate を \texttt{C=1} の前に置く順序が効いている。
  \item \textbf{\texttt{cores\_online = 4}、全ベンチ \texttt{agree = ok}}。
        SMPEN を立てたことで、\texttt{C=1} 下でもロックフリー・ジョブメールボックスの
        コヒーレンシが保たれている。逐次と並列の答えは全ベンチで一致した。
  \item \textbf{\texttt{sctlr = 0x00c5183d}}。D-cache OFF の \texttt{0x00c51839} に
        対しビット 2（C）が立ち、実機で D-cache が有効化されたことを裏付ける。
\end{itemize}

\subsubsection*{本命：メモリ律速はキャッシュでスケールする——ただし容量まで}

\begin{table}[htbp]
\centering
\caption{\bench{fill} の速度向上：D-cache OFF vs ON（8 パス、作業セット掃引）}
\label{tab:dcache-fill}
\small
\begin{tabular}{@{}rccl@{}}
\toprule
語数（作業セット） & OFF の速度向上 & ON の速度向上 & \\
\midrule
64          & \spd{0.38} & \spd{2.00} & \\
256         & \spd{0.61} & \spd{2.50} & \\
1{,}024 (4\,KB)   & \spd{1.13} & \spd{3.60} & \\
4{,}096 (16\,KB)  & \spd{1.27} & \spd{3.88} & \\
16{,}384 (64\,KB) & \spd{1.34} & \spd{4.01} & ← 完全スケール \\
65{,}536 (256\,KB)& \spd{1.24} & \spd{2.51} & ← 低下開始 \\
262{,}144 (1\,MB) & \spd{1.35} & \spd{1.40} & ← 帯域律速へ回帰 \\
\bottomrule
\end{tabular}
\end{table}

\noindent
表 \ref{tab:dcache-fill} が本実験の答えである。第 \ref{sec:rpi3} 節の結論
「メモリ律速はスケールしない」は、\textbf{D-cache OFF に固有の現象}であった。
D-cache ON では、作業セットがキャッシュに収まる限り\textbf{メモリ書込は
理想的にスケールする}（64\,KB で \spd{4.01}）。

しかしスケールするのは容量までである。作業セットが 256\,KB を超えると速度向上は
崩れ、1\,MB では \spd{1.40} まで戻る。これは\textbf{キャッシュ容量効果}である。
小さい \bench{fill} は同じ領域を 8 回書くため 2 回目以降はキャッシュに当たり、
実質\textbf{計算律速}になってスケールする。作業セットがキャッシュ（Pi 3 の
Cortex-A53 は L1 データ 32\,KB／コア＋共有 L2 512\,KB）を超えると、
書き戻しと充填で\textbf{再び帯域律速}に戻る。速度向上が 64\,KB で頂点、
256\,KB で崩れ始めるのは、この容量境界と整合する。

\subsubsection*{ウィンドウシステムの結論は覆るか——覆らない}

一見、これは「WM は並列化しない」（第 \ref{sec:rpi3} 節）を覆すように見える。
だが\textbf{覆らない}。理由が 2 つある。

第一に、\textbf{実フレームバッファはキャッシュ不可にマップしてある}
（第 \ref{sec:dcache-usb} 項、\texttt{mmu\_set\_range\_noncached}）。GPU が
スヌープしないため、これは D-cache ON でも変えられない。表 \ref{tab:dcache-fill}
の \bench{fill} は\textbf{通常のキャッシュ可能バッファ}での測定であり、
実描画はこの恩恵を受けない。

第二に、恩恵を受けるのは作業セットがキャッシュに収まる場合だけだが、
全画面ブリットは 1280$\times$800$\times$4 = 3.9\,MB でキャッシュを大きく超え、
仮にキャッシュ可能でも \spd{1.40} 程度に留まる。

したがって結論は次のように\textbf{精緻化}される：メモリ律速の処理は、
D-cache ON かつ\textbf{作業セットがキャッシュ内}なら 4 コアでスケールする。
だが実描画は（キャッシュ不可の FB へ、キャッシュを超える量を書くため）
どちらの条件も満たさず、依然スケールしない。中間バッファ上での
合成のように、キャッシュ内で完結する処理を挟めば話は別である。

\subsubsection*{副産物：D-cache は単一コア性能も大きく上げる}

計算律速側では、\bench{dining}(\spd{3.99})・\bench{primes}(\spd{2.05}) の
速度向上比は D-cache OFF とほぼ不変であった（元々スケールしていたため）。
一方で\textbf{絶対性能}は激変した。\bench{nqueens n=12} の 1 コア時間は
D-cache OFF の約 167{,}000\,\textmu s から\textbf{約 19{,}000\,\textmu s へ、
実に 8.8 倍高速化}している。クロック差（起動直後 2.33 倍）だけでは説明できず、
再帰探索が触る盤面配列 \texttt{cols[]} とスタックがキャッシュに載った効果である。
並列スケーリングとは独立に、D-cache は\textbf{計算律速の絶対性能そのもの}を
底上げする。

\subsubsection*{ただし本番投入はできない}

以上は\textbf{実験変種}（\texttt{DCACHE\_EXPERIMENT}、USB を落とした構成）での
結果である。本番カーネルに D-cache を入れるには、第 \ref{sec:dcache-usb} 項の
USB DMA 問題——ヒープとスタックのバッファへのバウンス、およびドライバ全域の
64 バイト境界化——を解く必要がある。それは本実験の範囲外であり、次段階の課題である。

\subsection{付録：実機シリアル出力（D-cache ON）}

参考のため、実機（コールドブート）から取得した起動時ハーネスの生出力を示す。

\begin{lstlisting}[language=,basicstyle=\ttfamily\scriptsize]
===== smpbench =====
cores_online = 4
D-cache      = ON   (SCTLR.C=1, SMPEN set)
sctlr        = 0x00c5183d

-- compute-bound (expect x3.2-x4.0) --
  bench           1-core us  N-core us  speedup
  dining n=5          51290      12828   x3.99  ok
  nqueens n=11         3716       1027   x3.61  ok
  nqueens n=12        19024       4865   x3.91  ok
  primes 200k         66761      32550   x2.05  ok

-- memory-store-bound: THE QUESTION (x1.0 = no benefit) --
  words           1-core us  N-core us  speedup
  fill 64                 2          1   x2.00  ok
  fill 256                5          2   x2.50  ok
  fill 1024              18          5   x3.60  ok
  fill 4096              70         18   x3.88  ok
  fill 16384            285         71   x4.01  ok
  fill 65536           1177        468   x2.51  ok
  fill 262144          4682       3330   x1.40  ok

-- pool round trip (the parallelise/don't break-even) --
  null                    0          1   x0.00  ok

===== end smpbench =====
\end{lstlisting}
\subsection{本番化の試みと限界：USB DMA は越えられなかった}
\label{sec:dcache-prod}

第 \ref{sec:dcache-usb} 項で「USB は構造的に難しい」と述べたが、実際に
本番化を試みた。結果は\textbf{部分成功}であり、その失敗の様相そのものが
価値ある結果なので記録する。

\subsubsection*{方式：全 USB 転送を非キャッシュ・バウンス経由に}

USB DMA の障害は、バッファがヒープ（\texttt{memget} は 8 バイト境界＝
64 バイトのキャッシュラインを隣と共有）やスタック上にあり、per-transfer の
clean/invalidate が\textbf{隣接データを巻き込む}ことだった。そこで
per-transfer 管理を避け、\textbf{全転送を専用の非キャッシュ・アリーナ経由に
バウンス}する方式を採った。DWC ドライバは元々、非境界バッファ用のバウンス
経路（\texttt{aligned\_bufs} + TX/RX コピー）を持つので、その切替を
「D-cache ON なら常にバウンス」に変え、アリーナを非キャッシュにマップした。
非キャッシュ領域は CPU も DMA も同一 RAM を直接見るので、原理上キャッシュ
管理命令が一切不要になる。

\subsubsection*{実機結果：起動と HDMI は成功、USB は全滅}

\begin{table}[htbp]
\centering
\caption{D-cache 本番化（フルカーネル）の実機結果}
\small
\begin{tabular}{@{}ll@{}}
\toprule
項目 & 結果 \\
\midrule
D-cache ON で起動 & $\checkmark$ brick なし（デスクトップ表示） \\
HDMI / フレームバッファ & $\checkmark$ 動作（FB の非キャッシュ再マップが有効） \\
set/way invalidate + SMPEN & $\checkmark$ 機能 \\
単一コア性能 & $\checkmark$ 8.8 倍（実験変種で実測） \\
\textbf{USB（マウス・キーボード・ethernet）} & $\times$ 全滅 \\
\bottomrule
\end{tabular}
\end{table}

\noindent
HDMI が正しく映ることは、非キャッシュ化の機構そのもの
（\texttt{mmu\_set\_range\_noncached}）が\textbf{正しく動作している証拠}である
（フレームバッファも同じ関数で非キャッシュ化しており、壊れていれば画面が
乱れる）。にもかかわらず USB は全滅した。Raspberry Pi 3 では
マウス・キーボード・ethernet がすべて LAN9514 ハブの背後にあるため、
USB が一つ壊れると全滅する。実機シリアルログには
\texttt{[webactor] autostart: ETH0 never came up} が記録された。

\subsubsection*{四つの修正を試みた——いずれも解決せず}

原因を、実機を最小限しか使わずオフラインのコードレビューで追った。

\begin{enumerate}[leftmargin=1.6em]
  \item \textbf{DSB バリア追加。} 非キャッシュ書き込みもライトバッファに
        滞留し得るため、DMA 起動前に DSB で RAM へ排出。\textbf{必要だが
        不十分}——単独では USB は復活しなかった。

  \item \textbf{\texttt{packet\_count} オーバーフロー修正。} \texttt{packet\_count}
        は 10 ビット（最大 1023）で、バウンスアリーナを 128\,KB に拡大した際に、
        小さい \texttt{max\_packet\_size} で桁溢れし得ることを発見。しかし
        \textbf{赤鰊だった}——スプリット転送では \texttt{transfer.size} が
        \texttt{max\_packet\_size} に事前制限され \texttt{packet\_count} は
        常に 1 で溢れない。撤回した。

  \item \textbf{スプリット転送ロジックの精査。} Pi 3 は全転送がスプリット
        （ハブ背後）なので、バウンスとスプリット再試行の相互作用を徹底
        トレースした。\textbf{ロジックは正しかった}——\texttt{dma\_address} の
        保持、再アーム、コピーバックのオフセットはすべて整合していた。

  \item \textbf{アリーナ非キャッシュ化の SMP タイミング修正。} 最有力仮説。
        非キャッシュ化を実行時（\texttt{usbinit} 内）に、しかも現在コアだけ
        TLB 無効化で行っていたが、SMP のワーカーコアはそれ\textbf{以前}に
        アリーナがキャッシュ可能なまま起動しており、コア間でマッピングが
        不整合だった。アリーナのアドレスはコンパイル時に確定するので、
        \texttt{mmu\_init} で全コア起動前・D-cache 有効化前に非キャッシュを
        ベイクするよう変更した。しかし\textbf{これも USB を復活させなかった}。
\end{enumerate}

\subsubsection*{結論：本番化は USB DMA に阻まれた}

四つの修正を尽くしても USB は動かなかった。バウンスロジックとスプリット
処理は精査済みで正しく、コヒーレンシ機構も（HDMI が証明する通り）機能して
いる。それでも USB が全滅する——真因は、これらより深い、実機レベルの精密
診断（USB 列挙の詳細なシリアルログ等）を要する相互作用にある。手元の
シリアルアダプタ（Prolific、macOS で不安定）ではそこまで追えなかった。

したがって本レポートの範囲では、\textbf{D-cache は計算と表示には本番投入
できるが、USB DMA が最後の、そして越えられない壁である}と結論する。第
\ref{sec:dcache-usb} 項の構造的予測——USB のバッファ出自が本質的な障害——が、
バウンスという迂回策すら退ける形で、より強く裏付けられた。8.8 倍の単一コア
高速化という果実は実在するが、それを本番で収穫するには、USB ドライバの
DMA 経路をより根本的に作り直す必要がある。これは明確に次段階の課題である。
